\section{A sound three-residue encoding by tagging one atom operand}
\label{sec:tagged-three}

The existing coefficient-three counterexamples concern constant-offset
encodings of equality.  The following encoding leaves implication and
universal quantification at coefficient three and places the atom tag in
one Cantor operand instead.  It has an exact proof-theoretic equivalence
with the original four-tag numerical calculus.  This is a mathematical
construction; it is not by itself a claim of a smaller literal machine.

Write $p(a,b)=(a+b)(a+b+1)/2+b$.  Denote the old constructors by
\[
 I_4(A,B)=4p(A,B),\quad M_4(a,b)=4p(a,b)+1,\quad
 U_4(a,B)=4p(B,a)+2,\quad E_4(a,b)=4p(a,b)+3.
\]
The proposed constructors are
\[
\boxed{\begin{aligned}
 I_3(A,B)&=3p(A,B),& U_3(a,B)&=3p(B,a)+2,\\
 M_3(a,b)&=3p(a,2b)+1,& E_3(a,b)&=3p(a,2b+1)+1.
\end{aligned}}
\tag{T3}\label{eq:tagged-three-constructors}
\]
Keep numerical truth $0$, numerical contradiction $1$, the six ordinary
head expressions, the separate diagonal head, the guarded gate, the
restricted propositional head, MP, and theorem generalization, replacing
only the four constructors by \eqref{eq:tagged-three-constructors}.
Negation remains $N_i(A)=I_i(A,1)$.

\subsection{Total decoding and the exceptional small codes}

For a positive new code $n$, residue zero modulo three means implication,
and residue two means universal quantification.  For residue one, unpair
$(n-1)/3$ as $(a,s)$.  Even $s=2b$ gives $M_3(a,b)$; odd $s=2b+1$ gives
$E_3(a,b)$.  The three residue classes and these two parity cases are
exhaustive and disjoint.  Thus there are no unused or ambiguous codes.
In particular, equality is not an implication as in the rejected
constant-offset construction.

As before, $I_i(0,0)=0$ is the sole zero alias.  Decode zero as truth.
Also $M_i(0,0)=1$ in both encodings, so the designated contradiction is
literally unchanged.  Atoms have no recursive formula children, but their
numerical label arguments are still smaller than their code.  Every
constructor argument of every nonzero constructor value is strictly
smaller than that value; for universal quantification this includes both
the body code and the binder label.  For example
$b\le2b\le p(a,2b)<3p(a,2b)+1$ except at the separately immediate
zero case.  The analogous inequalities for equality are stronger.
These inequalities justify the recursions below in iPRA.

There is a useful additional freshness fact.  Every label occurring
anywhere in the formula decoded from $B$ is strictly below $B$, with code
zero handled separately.  This follows by bounded structural induction
from the same constructor inequalities.  Consequently the new diagonal
clause $B\to U_3(B,B)$ is a valid fresh-label vacuity clause, just as in
the old encoding.  This fact is compatible with the direct proof
translation below but is not being used as a substitute for it.

\subsection{Recode formula codes and variable labels by the same permutation}

A recoding that fixes variable labels would disturb the diagonal and gate
interfaces, because those interfaces use a formula code as a variable
index.  Instead define a single function $f:\mathbb N\to\mathbb N$ that
acts on both kinds of numerical arguments.  Put $f(0)=0$ and require
\[
\begin{aligned}
 f(I_4(a,b))&=I_3(f(a),f(b)),&
 f(M_4(a,b))&=M_3(f(a),f(b)),\\
 f(U_4(a,b))&=U_3(f(a),f(b)),&
 f(E_4(a,b))&=E_3(f(a),f(b)).
\end{aligned}
\tag{F}\label{eq:tagged-three-f}
\]
For nonzero input the old residue and Cantor projections choose exactly
one clause.  Its arguments are smaller than the input.  The only
remaining instance $I_4(0,0)=0$ satisfies the equation automatically.

Define $g$ in the opposite direction using the new decoding just given,
with $g(0)=0$, replacing the subscript 3 on the left of each equation in
\eqref{eq:tagged-three-f} by 4 on the right.  These are total
primitive-recursive functions.  For an explicit implementation of the
course-of-values recursion, successively construct the finite value table
through the input code; each new cell reads only earlier cells.  The
magnitude of the newly computed values need not be bounded by the input.
The table construction itself is a primitive recursion and supplies its
own code and bound.

Bounded strong induction, expressed as a Boolean conjunction over the
prefix of the input range, proves
\[
 g(f(n))=n,\qquad f(g(n))=n.
\tag{Inv}\label{eq:tagged-three-inverse}
\]
Indeed, the outer constructors decode uniquely; their argument equations
reduce the assertion to its already proved smaller cases.  The equations
are quantifier free in the language with the displayed PR functions, so
this proof uses only iPRA.  Thus $f,g$ are mutually inverse permutations,
and $f$ is injective.  Their first values are
\[
 (f(0),\ldots,f(8))=(0,1,2,7,3,4,5,13,6).
\]
In particular, both zero and one are fixed.  The permutations are
\emph{not} monotone: $3<4$ but $f(3)=7>3=f(4)$.

\subsection{Exact seed and guard correspondence}

Each of the six ordinary seeds is a finite expression in $I,U,M,E$, the
constant $1$, and its displayed parameters.  For clarity their signatures
are $h_1(D,H)$, $h_2(D,A,H)$, $h_3(D,A,H)$, $h_4(D,A,H)$,
$h_5(D,H)$, and $h_6(D,A,H)$.  All binder and atom labels in them are
supplied parameters.  There is no distinguished variable label $3$ or
other fixed nonzero label hidden in these seeds.  The source increments
by 2 and 3 are old constructor tags, not variable-name constants.
Explicitly, with $\exists_d X=N(U(d,N(X)))$, the target seeds are
\[
\begin{aligned}
h_1(D,H)&=(\exists_D U(D,H))\to H,\\
h_2(D,A,H)&=U(D,U(A,H))\to U(A,U(D,H)),\\
h_3(D,A,H)&=M(H,D)\to(E(H,A)\to M(A,D)),\\
h_4(D,A,H)&=E(H,D)\to(M(A,H)\to M(A,D)),\\
h_5(D,H)&=\exists_D E(H,D),\\
h_6(D,A,H)&=U(D,A)\to(U(D,A\to H)\to U(D,H)).
\end{aligned}
\]
Here every arrow means $I_3$, and every constructor is exactly (T3).
The Q variant exchanges the two antecedents of $h_6$.
It follows by finite congruence from \eqref{eq:tagged-three-f} that
\[
 f(h_j^{(4)}(\vec a))=h_j^{(3)}(f(\vec a))
\]
for every one of the six seeds, including every zero-parameter case.
The exchanged sixth head used by Q variants obeys the same equation.
Negation, conjunction, and existential abbreviation commute with $f$
because $f(1)=1$.

The diagonal maps \emph{exactly}, despite its use of its own formula code
as a binder:
\[
 f(I_4(c,U_4(c,c)))=I_3(f(c),U_3(f(c),f(c))).
\]
For the gate use its original definitions
\[
\begin{aligned}
 V&=I(M(a,b),U(c,M(a,b))), &P&=I(c,M(c,b)),\\
 S&=U(b,(M(a,b)\to P)\wedge(M(c,b)\to c)),&
 Q&=U(c,S\to N(P)),\\
 G(a,b,c)&=N(I(V,Q)).
\end{aligned}
\]
Every occurrence of $c$, whether as a formula, atom label, or binder, is
mapped to the same number $f(c)$.  Hence
\[
 f(G_4(a,b,c))=G_3(f(a),f(b),f(c)).
\]
The actual guard is preserved in both directions:
\[
 c\ne a,b\quad\longleftrightarrow\quad f(c)\ne f(a),f(b),
\]
by \eqref{eq:tagged-three-inverse}.  A failed gate is zero in both
systems.  No stronger guard, restricted variable class, or unproved
renaming rule is needed.

\subsection{The numerical domination guard is handled by proof expansion}

For either $I_i$ let
\[
 \rho_i(D,A,H,u)=
 I_i(I_i(I_i(H,u),D),I_i(I_i(D,H),I_i(A,H))).
\]
The axiom row has the numerical guard $u\le D$.  It would be incorrect
to assume that this guard is preserved by $f$ or $g$; the small example
above already disproves that assumption.

Instead use the finite generic restricted-implication derivation
\path{M6.Logic.RestrictedImplication.rho8All}.
Its \texttt{Rules} interface requires MP, the restricted row, reflexivity
of the comparison, and propagation of a bound into either argument of
implication.  For both $i=3,4$ the last requirements follow from
\[
 a,b\le p(a,b)\le I_i(a,b).
\]
The retained source explicitly derives $K$, Peirce, composition and
exchange, and then the unrestricted $\rho_i(D,A,H,u)$ for arbitrary four
parameters.  These are finite proof templates with only primitive
constructor operations and bound proofs through implication children.
Their expansion is a primitive-recursive function of the four parameters,
with a fixed finite bound on its number of rows.  Their checking equations
and the displayed arithmetic bounds are provable by quantifier-free
induction.  There is no dependence here on four-residue parsing or on the
monotonicity of the recoding.

Thus an old guarded-rho leaf is translated to a target \emph{derivation}
of its recoded unrestricted instance.  The reverse translation uses the
same generic template for coefficient four.  This single expansion is the
only place where the recoding is not simply one primitive axiom row to one
primitive axiom row.  A copy of the cited generic source is retained as
\texttt{\detokenize{reference/RestrictedImplication.lean}}; its SHA-256 is
\begin{quote}\small\ttfamily
00af59bb2f970f780b834b721cccd06c4ed5b5ee13fbfddb0deb0139b4c62581.
\end{quote}

The archive additionally contains a complete formal expansion,
\texttt{rho\_generic\_template.json.gz}, and an independent checker that
does not read the Lean source.  The expansion uses four unrestricted
formal variables, 32,455 primitive guarded-rho rows and 35,813 MP rows,
for 68,268 proof rows in total.  The checker verifies that every
primitive guard is justified by syntactic subtree inclusion of the
fourth parameter in the first.  Each such inclusion is interpreted as
a numerical inequality by reflexivity and the two implication bounds.
Every MP match is literal term equality, so all matches survive every
valuation of the four variables.  Thus this is a fixed universal proof
template, not a numerical sample.  Its compressed SHA-256 is
\begin{quote}\small\ttfamily
9fd2d5082a97dc051e4b56cb6e705f1d6f4f9bcabe82eff7fcbb954c6d285278.
\end{quote}
The separate concrete certificate checks the nonmonotone example
$D=3,u=7$ at coefficient three with all primitive guards valid even
though the final unrestricted instance fails its direct guard.

\subsection{Primitive-recursive proof compilers in both directions}

Encode derivations as finite lists whose MP and Gen premises refer only
to earlier rows, with primitive-recursive row checking.  Define $F$ by a
single scan of an old derivation.  Maintain a list of translated rows and
an index map recording the translated last row of every processed source
row.  Recode head parameters by $f$, replace a rho leaf by the fixed
expansion just described, map an MP row to MP of its two recorded indices,
and map a Gen row at binder $d$ to Gen at binder $f(d)$.  Use an arbitrary
fixed default output on invalid input.  Appending lists, relocating indices,
and all numerical parameter computations are PR; hence $F$ is PR.
Define $G$ symmetrically with $g$.

The induction invariant is a bounded Boolean check that all processed
source rows have valid target proofs, their recorded indices are in range,
and the formula at each recorded index is the recoding of the corresponding
source formula.  The seed, guard, and rho cases are the preceding equations.
The MP case follows from implication homomorphism, including
$I_i(0,0)=0$.  The Gen case follows from universal homomorphism.  This
quantifier-free prefix induction proves in iPRA
\[
\begin{aligned}
 \operatorname{Prf}_{\mathcal C_4}(d,A)&\longrightarrow
       \operatorname{Prf}_{\mathcal C_3^{\rm tag}}(F(d),f(A)),\\
 \operatorname{Prf}_{\mathcal C_3^{\rm tag}}(e,B)&\longrightarrow
       \operatorname{Prf}_{\mathcal C_4}(G(e),g(B)).
\end{aligned}
\]
Because $f(1)=g(1)=1$, the conclusion is the exact refutation equivalence
\[
\boxed{\mathsf{iPRA}\vdash
 (\exists d\,\operatorname{Prf}_{\mathcal C_4}(d,1))
 \longleftrightarrow
 (\exists e\,\operatorname{Prf}_{\mathcal C_3^{\rm tag}}(e,1)).}
\]
The same statement holds for the Q-oriented sixth head.  Composing with
the pinned written calibration of $\mathcal C_4$ requires no new
PA-to-relational or Foundation proof compiler.

\subsection{What an implementation must actually change}

With the stack constructor consuming its first operand from the stack and
its second operand from visible register $T$, an implication uses the
coefficient-three constructor unchanged.  For membership, double $T$
\emph{before} pairing and add one after construction.  For equality,
double $T$, add one \emph{before} pairing, construct with coefficient
three, and add one afterwards.  Universal quantification remains the
coefficient-three constructor followed by two increments in the old
body/binder operand order.  The doubling macro must preserve the saved
stack, every logical parameter register, and its stated scratch invariants.

The abstract proof-search compiler must be checked with these exact
constructor functions.  Its standard residue tests remain valid: all
implications have residue zero and all universal codes residue two, so
those emitted heads are never the contradiction 1.  Zero implication
still forces both operands zero.  The dominated-rho selector still
chooses $D=u+x$ in the \emph{new} numerical coding; the proof translation
does not assert that $f(u+x)=f(u)+f(x)$.  The scheduler, command tags,
Cantor stack coding, and initial candidate counter are separate data
structures.  Their correctness follows from their own target command
contracts and exhaustive scheduling, not from a claimed homomorphism
of entire old and new register executions.

The supplied \texttt{recoding.py} independently checks 10,000 inverse
instances in each direction, 4,608 constructor identities, all six-head
and Q-head identities on small parameters, guarded gate identities,
diagonal identities, and unrestricted rho-expression identities.  These
finite checks support the displayed universal proof and expose the
nonmonotone guard example.  They do not replace the written induction
argument or claim a newly built proof-assistant endpoint.  The state count
of a compiled implementation must be determined separately from this
proof-theoretic result.

\subsection{A further variant with a shared first-operand transformation}

A related encoding may be worth pricing separately:
\[
\begin{aligned}
 I_2(A,B)&=2p(2A,B),& U_2(a,B)&=2p(2B+1,a),\\
 M_2(a,b)&=2p(2a,b)+1,& E_2(a,b)&=2p(2a+1,b)+1.
\end{aligned}
\]
Given $n$, unpair $\lfloor n/2\rfloor$ as $(c,b)$.  For even $n$,
even $c=2a$ denotes $I_2(a,b)$ and odd $c=2a+1$ denotes $U_2(b,a)$.
For odd $n$, those two cases denote $M_2(a,b)$ and $E_2(a,b)$.
The four images are disjoint, cover all codes, preserve zero and one,
and have the same strict-argument bounds.  The identical homomorphic
recursion on both codes and labels consequently gives inverse PR
permutations with the old encoding.  The rho expansion remains valid
because $A,B\le I_2(A,B)$.

This variant places a doubling operation on the first Cantor operand for
\emph{every} constructor, with one additional increment for universal and
equality.  That uniformity suggests a possible shared-pairing
implementation, whereas the three-residue proposal doubles only atom
operands.  No validated physical implementation of this variant is
established here; the separate pricing prototype fails a source contract,
and no state saving is asserted.  A command compiler
would also replace old residue arguments by exact constructor-image
separation: implication and universal codes are now both even, separated
by the parity of their first decoded Cantor component.  The supplemental
\texttt{tagged2\_check.py} checks the inverse and seed identities on finite
samples.  It is an implementation aid for further research, not a
completion of that source-engineering task.
